\documentclass[10pt,a4paper]{amsart}
\usepackage[margin=19mm]{geometry}
\usepackage{amsmath,amssymb,mathtools}
\usepackage[scr=rsfs,cal=euler]{mathalfa}
\usepackage{xfrac}
\usepackage[unicode,hidelinks,bookmarks=false]{hyperref}
\newcommand{\ie}{i.e.\ }
\newcommand{\cf}{cf.\ }
\newcommand{\resp}{respectively\ }
\newcommand{\Subsection}{\subsection}
\providecommand{\nobreakdash}{\nobreak\hbox{-}}
\setlength{\emergencystretch}{2em}
\multlinegap0pt
\usepackage{tikz}
\usetikzlibrary{cd}
\usepackage{amssymb}
\usepackage{xcolor}
\usepackage[shortlabels]{enumitem}
\usepackage{mathtools}
\newtheorem{theorem}{Theorem}
\newtheorem{lemma}{Lemma}
\newtheorem{proposition}{Proposition}
\newcommand{\C}{\mathbb{C}}
\newcommand{\JS}[1]{\textcolor{teal}{[JS: \footnotesize #1]}}
\newcommand{\Stab}{\mathrm{Stab}}
\newcommand{\cS}[1]{\textcolor{purple}{[SN: \footnotesize $\blacksquare$ #1 $\blacksquare$]}}
\newcommand{\dil}{\mathrm{Dil}}
\newcommand{\excluster}[1]{\widetilde{\mathbf{#1}}}
\newcommand{\exmatrix}[1]{\widetilde{#1}}
\newcommand{\mut}{\mathrm{mut}}
\newcommand{\seed}{\Sigma}
\newcommand{\seeds}{\mathsf{Seeds}}
\newcommand{\var}{\mathcal{A}}
\title{Mysterious points in keys but not trees}
\author{Scott Neville, Jos\'e Simental}
\date{Source: arXiv:2603.17890v1 (CC BY 4.0)}
\begin{document}
\maketitle

\noindent\emph{Source and licence.} Mysterious points in keys but not trees. Version arXiv:2603.17890v1 (CC BY 4.0), distributed by arXiv under the Creative Commons Attribution 4.0 International licence, http://creativecommons.org/licenses/by/4.0/, as recorded in the arXiv licence metadata for this exact version and shown on the abstract page https://arxiv.org/abs/2603.17890v1. Full licence notice: \texttt{licenses/arxiv-2603.17890-CC-BY-4.0.txt}.\par\smallskip

\noindent\textit{Editorial scope.} Counted unit: the article's proposition prop:3-vertices (source0.tex lines 976--991). The article's complete original proof of it is delivered with it (source0.tex lines 992--1009, one \texttt{proof} environment of 18 lines). No mathematics is written, repaired or re-derived: the statement and the proof are the article's own lines, in the article's own notation, and no retained line is substituted or re-flowed.  Retained uncounted as the source-local context the proof needs: lines 166--174; lines 386--403; lines 413--415; lines 600--602; lines 667--671.  Nothing was omitted from any retained span, so the package carries no omission record.  No retained line contains a citation, so the wrapper carries no bibliography and no bibliography entry.  The retained lines contain the article's own diagram code, which the wrapper typesets with the standard tikz packages; no image file is delivered and the package declares no asset dependency.  Editorial formatting only, in addition: an AMS \texttt{amsart} preamble that restates the article's own declarations and macros used by the retained lines, namely the environments \texttt{theorem}, \texttt{lemma}, \texttt{proposition} on separate counters, together with the standard packages those lines need.  The retained environments print in this document as Theorem 1, Lemma 1, Proposition 1 in order of appearance, whereas the article prints the same items with the numbering of its own full document.  The complete native source of the article as distributed in the arXiv e-print for this exact version is bundled unchanged as \texttt{sources/196658/source0.tex}.
\begin{theorem}\label{thm:main2-intro}
Consider the quiver
\[\begin{tikzcd}
	3 && 1 && 2
	\arrow["a", from=1-1, to=1-3] %a=2
	\arrow["b"', from=1-5, to=1-3] %b=3
\end{tikzcd}\]
Assume that $\gcd(a,b)=1$, $\min(a,b) \geq 2$. Then, the corresponding cluster algebra has mysterious points. 
\end{theorem}

\begin{lemma}\label{lem:creating-deep-points-frozens}
    Let $\seed$ be a seed such that $Q^{\mut}$ is acyclic, and let $\widetilde{\seed}$ be obtained from $\seed$ by adding $\ell$ arbitrary frozen rows to the exchange matrix. Let 
    \[
    %p = (p_1, \dots, p_n, p'_1, \dots, p'_n, p_{n+1}, \dots, p_{n+m}) \in \var(\seed)
    p = (p_1,p'_1,\dots, p_n, p'_n, p_{n+1}, \dots, p_{n+m}) \in \var(\seed)
    \]
    and define the point
    \[
    %\widetilde{p} = (p_1, \dots, p_n, p'_1, \dots, p'_n, p_{n+1}, \dots, p_{n+m}, 1, \dots, 1) \in \C^{2n} \times (\C^{\times})^{m+\ell}.
    \widetilde{p} = (p_1, p'_1,\dots, p_n, p'_n, p_{n+1}, \dots, p_{n+m}, 1, \dots, 1) \in \C^{2n} \times (\C^{\times})^{m+\ell}.
    \]
    Then
    \begin{enumerate}
        \item $\widetilde{p} \in \var(\widetilde{\seed})$.
        \item $p$ belongs to a cluster torus in $\var(\seed)$ if and only if $\widetilde{p}$ belongs to a cluster torus in $\var(\widetilde{\seed})$.
        \item The groups $\Stab_{\dil(\seed)}(p)$ and $\Stab_{\dil(\widetilde{\seed})}(\widetilde{p})$ are isomorphic.
    \end{enumerate}
\end{lemma}

\begin{theorem}\label{thm:main-notintro}
Let $\seed = (\exmatrix{B}(Q), \excluster{x})$ be a seed such that $Q^{\mut}$ is a tree. Then, $\var(\seed)$ has no mysterious points. 
\end{theorem}

\begin{lemma}\label{lem:rank2}
Rank $2$ cluster algebras have no mysterious points.
\end{lemma}

\begin{proposition}\label{prop:so-may-deep}
For a seed $\seed  = (\exmatrix{B}, \exmatrix{x})$, suppose that for all $\seed' = (\exmatrix{B}', \exmatrix{z}) \in \seeds(\seed)$, we have $|b_{1,i}'| \geq 2$ for $i$ in $2, \dots, n$. 
Fix $p\in \var(\seed)$ such that $x_1(p) = x'_1(p)=0$ and $x_j(p) \neq 0$ for $j\neq1$. %(where $x'_j$ is the new cluster variable in $\mu_j(\seed)$).
Then $p$ is deep. 
\end{proposition}

\begin{proposition}\label{prop:3-vertices}
Consider the cluster algebra $A$ with exchange matrix associated to the quiver
\[\begin{tikzcd}
	{\boxed{\overline{3}}} && {\boxed{\overline{1}}} && {\boxed{\overline{2}}} \\
	3 && 1 && 2
	\arrow[from=1-3, to=2-3]
	\arrow[from=2-1, to=1-1]
	\arrow["a", from=2-1, to=2-3]
	\arrow[from=2-5, to=1-5]
	\arrow["b"', from=2-5, to=2-3]
\end{tikzcd}\]
\begin{enumerate}
    \item If either $\min(a,b) = 1$ or $\gcd(a,b) > 1$, then $A$ has no mysterious points.
    \item Else, $A$ has mysterious points. 
\end{enumerate}
\end{proposition}

\begin{proof}
We assume that $a$ and $b$ are not both equal to $1$; otherwise the result follows from Theorem \ref{thm:main-notintro}. The cluster dilation group is
\[
\dil(\seed) = \{t = (t_1, t_2, t_3, t_{\overline{1}}, t_{\overline{2}}, t_{\overline{3}}) \in  (\C^{\times})^{6} \mid t_1^{a}t_{\overline{3}} = t_3^{a}t_2^{b}t_{\overline{1}} = t_{1}^{b}t_{\overline{2}} = 1\}.
\]
Let $p \in \var$ have trivial stabilizer in $\dil(\seed)$. If $p \in \mathcal{O}_{\emptyset} = T(\seed)$ we are done. If $p \in \mathcal{O}_{\{3\}}$ then $x_1(p)x_2(p)\neq 0$, so we may freeze vertex $2$ in order to reduce to the case of Lemma \ref{lem:rank2}. Similarly if $p \in \mathcal{O}_{\{1\}}$. 

It remains to treat the case of the maximal independent sets $\{2,3\}$ and $\{1\}$. If $p \in \mathcal{O}_{\{2,3\}}$ but either $x'_2(p)$ or $x_3'(p) \neq 0$, we mutate at the corresponding vertex to reduce to the case of the previous paragraph. So we assume $x_2(p) = x'_2(p) = x_3(p) = x_3'(p) = 0$. If, say, $a > 1$, this implies that $(\zeta, 1, 1, 1, 1, 1) \in \Stab_{\dil(\seed)}(p)$, where $\zeta$ is a primitive $a$-th root of unity. So this case cannot happen.

Now assume that $p \in \mathcal{O}_{\{1\}}$. We may also assume that $x'_1(p) = 0$. If $\gcd(a,b) > 1$, we have $(1, \xi, 1, 1, 1, 1) \in \Stab_{\dil(\seed)}(p)$, where $\xi$ is a primitive $\gcd(a,b)$-root of unity, a contradiction. If $\min\{a,b\} = 1$, say $a = 1$, then we apply a strategy similar to that of the proof of Case 2.1 of Theorem \ref{thm:main-notintro}, and we see that $p \in T(\mu_1\mu_3(\seed))$, indeed
\[
x''_1 = \frac{x'_3 + x_2^{b}x_{\overline{1}}}{x_1} = 
\frac{\frac{1 + x_{\overline{3}}x_1}{x_3}+ x_2^{b}x_{\overline{1}}}{x_1} = 
\frac{1 + x_{\overline{3}}x_1 + x_{3}x_2^{b}x_{\overline{1}}}{x_1x_3} = 
\frac{x'_1 + x_{\overline{3}}}{x_3}
\]
so $x''_1(p) \neq 0$ and thus $p$ is not mysterious. We are done with (1). Statement (2) is Theorem \ref{thm:main2-intro}. %For statement (2), we may use Lemma \ref{lem:creating-deep-points-frozens} to eliminate the frozen variables. Since in this case $\gcd(a,b) = 1$, any point in the nonempty locus $\{x_1 = x_1' = 0\}$ will have trivial stabilizer, but it will also be deep by Proposition \ref{prop:so-may-deep} \JS{does the proposition apply here?} \cS{yes, and I think the rewrite makes it clear.}. So any such point is mysterious. 
\end{proof}

\end{document}
